% !PRINTED 464
\weberpage{464}
%<<CONT
\weq{10}{
\begin{aligned}
&\beta\partial = Ax,\\
&\alpha c - \gamma a = Cx,\\
&\beta c - \delta a + \alpha\partial = Bx,\\
&\beta c - \delta a - \alpha\partial = y,\\
&\beta c - \delta a = \frac{Bx+y}{2}, \quad \alpha\partial = \frac{Bx-y}{2},\\
&4n = y^2 - Dx^2.
\end{aligned}
}

If we assume $A$ and $n$ to have no common divisor, then $\partial = 1$, $a = n$ must hold; for a prime dividing $x$ and $\partial$ would have to divide $y$ as well, and consequently $\alpha c - \gamma a$ and $\beta c - \delta a$, and, since $\alpha\delta - \beta\gamma = 1$, also $a$ and $c$, whereas $a$, $c$, $\partial$ are to have no common divisor. From (10), $\alpha$, $\beta$, $\gamma$, $\delta$, $c$ are determined in the same way as in the preceding section. There result first of all the two congruences
\[
cAx \equiv \frac{Bx+y}{2}, \quad c\,\frac{Bx-y}{2} \equiv Cx \ (\mathrm{mod}\ n),
\]
of which the one follows from the other if $x$ is coprime to $n$, and the condition
\[
c \equiv 0 \ (\mathrm{mod}\ 2)
\]
is compatible with them. Further it follows that:
\weq{11}{
\begin{aligned}
&\alpha = \frac{Bx-y}{2}, \quad n\gamma = -Cx + c\,\frac{Bx-y}{2},\\
&\beta = Ax, \qquad\qquad n\delta = -\frac{Bx+y}{2} + Acx.
\end{aligned}
}

If now $D$, and hence $B$, is odd, then, since by (9) $\alpha$ and $\delta$ are either both even or both odd, $x$ and $y$ must be even; accordingly $\beta$ and $\gamma$ are even, and $\alpha$ and $\delta$ must be odd. Consequently, of the two numbers $\frac{1}{2}x$, $\frac{1}{2}y$ one must be even and one odd. But this is independent of how $A$ and $C$ behave with respect to the modulus 2.

If, therefore, one of the three quantities (2) satisfies the equation (7), then so do the other two, and we cannot separate these three from one another.

We must therefore assume $B$ to be even, and for simplification we set:
\weq{12}{D = -4m,}
in which $m$ is a positive integer.

We distinguish:
%CONT>>
